\chapter{Environment and deployment}
\label{ch:ch2}
\section{Requirements and version boundaries}
The deployment requires a disposable Linux Kubernetes cluster, a compatible Cilium
installation, Helm and an image build tool. Select and record a supported combination
of Kubernetes, Cilium, kernel and container runtime. Do not replace an existing
cluster's network plugin merely to run this project.

The committed Falco configuration targets chart \textbf{9.1.0} and Falco
\textbf{0.44.1}, using the modern-eBPF driver with kernel BTF support. The historical
kernel-module configuration is not part of this deployment procedure. Falco and
Cilium have distinct roles: syscall-based detection and network-policy enforcement,
respectively.

Direct Python dependencies and the chart version are specified, but container base
layers and transitive dependencies are not fully locked. The implementation therefore
does not claim bit-for-bit reproducibility. Record image digests and ruleset versions
when collecting integration evidence.

All commands below assume the repository root as the working directory. The image
import example applies to a single-node containerd environment; other runtimes need
their corresponding image-loading mechanism.

\section{Building and deploying the target}
\begin{lstlisting}[language=bash]
kubectl config current-context
docker build -t vulnapp:lab-v2 material/vulnapp
docker build -t falco-handler:lab-v2 material/falco_handler/src
docker save vulnapp:lab-v2 falco-handler:lab-v2 | \
  sudo ctr -n k8s.io images import -
# Establish quarantine policy before enabling automatic response.
kubectl apply -f material/manifests/isolate-tainted-policy.yaml
kubectl apply -f material/manifests/vulnapp-deployment.yaml
kubectl apply -f material/manifests/vulnapp-service.yaml
kubectl apply -f material/manifests/vulnapp-pdb.yaml
kubectl -n default rollout status deployment/vulnapp
\end{lstlisting}
The target has two replicas and a readiness probe. A NodePort Service on port 32080
selects pods labelled \texttt{app=vulnapp}. The ServiceAccount token is deliberately
mounted for the credential-exposure scenario. This choice must not be copied into
an application that does not require Kubernetes API access.

\section{Deploying the authenticated handler}
Generate one shared token and keep it available in the same shell until Falco has
been configured. Never commit a real token to the repository.
\begin{lstlisting}[language=bash]
export FALCO_WEBHOOK_TOKEN="$(openssl rand -hex 32)"
kubectl -n default create secret generic falco-handler-auth \
  --from-literal=webhook-token="$FALCO_WEBHOOK_TOKEN" \
  --dry-run=client -o yaml | kubectl apply -f -
kubectl apply -f material/falco_handler/falco-handler-rbac.yaml
kubectl apply -f material/falco_handler/falco-handler-service.yaml
kubectl apply -f material/falco_handler/falco-handler-networkpolicy.yaml
kubectl apply -f material/falco_handler/falco-handler-deployment.yaml
kubectl -n default rollout restart deployment/falco-handler
kubectl -n default rollout status deployment/falco-handler
\end{lstlisting}
The restart also refreshes Secret-backed environment variables when rotating the
token on an existing installation. The handler Service uses port 80, forwarding to
container port 5000. Its ingress policy selects only Falcosidekick pods in the
\texttt{falco} namespace. The namespaced Role grants pod get/list/patch/delete;
workload filtering is enforced by the handler, not by label-scoped RBAC.

\section{Configuring the event pipeline}
Create the Falco namespace and restrict the upstream receiver before installing
Falcosidekick. The committed values use pod networking; verify the rendered labels
and network behaviour in the selected environment.
\begin{lstlisting}[language=bash]
kubectl create namespace falco --dry-run=client -o yaml | kubectl apply -f -
kubectl apply -f material/falco/falcosidekick-networkpolicy.yaml
helm repo add falcosecurity https://falcosecurity.github.io/charts
helm repo update
(
  umask 077
  FALCO_VALUES_DIR=$(mktemp -d)
  trap 'rm -f "$FALCO_VALUES_DIR/values.yaml"; rmdir "$FALCO_VALUES_DIR"' EXIT
  cp material/falco/falco-values.example.yaml "$FALCO_VALUES_DIR/values.yaml"
  perl -pi -e 's/CHANGE_ME/$ENV{FALCO_WEBHOOK_TOKEN}/g' \
    "$FALCO_VALUES_DIR/values.yaml"
  helm upgrade --install falco falcosecurity/falco --version 9.1.0 \
    --reset-values -n falco -f "$FALCO_VALUES_DIR/values.yaml" --wait
)
unset FALCO_WEBHOOK_TOKEN
\end{lstlisting}
Review intentional local overrides before using \texttt{--reset-values}. Reusing
old values can retain duplicate custom-rule definitions. Helm release metadata can
contain supplied configuration secrets, so access to those objects must be restricted.
The shared token authenticates a request; unencrypted HTTP does not provide transport
confidentiality. Production deployment would require stronger transport and secret controls.

\section{Migration and operational checks}
An existing ClusterRoleBinding is not removed by applying a RoleBinding. Inspect
obsolete global bindings and their consumers before deleting any of them. Confirm
that the handler cannot patch pods outside the intended namespace. The repository
README lists the legacy object names and current deployment commands.

Before enabling response, check both directions of the event path: Falco must reach
Falcosidekick and Falcosidekick must reach the handler, while unrelated pods must not.
NetworkPolicies are additive, so inspect other policies and administrator/node access.
For detection-only observations, deploy Falco without enabling the response handler;
otherwise an accepted event can interrupt the scenario by quarantining the target.
